The sequence diagrams below show the four flows whose behaviour is spread over
several components: a break from start to automatic end, signing in with a
token refresh, the video limit of the TV, and recording a daily game. Message
names follow the code (stores, API endpoints and in-app browser events).

\subsection{Start a Break and Automatic End}
\label{sec:seq-break}

The session store is the single source of truth for the countdown. It stores
the planned end time, so a tick after the app returns from the background
immediately shows the correct remaining time (FR 2.4). When the countdown
reaches zero, the store stops the media and saves the record before the
summary is shown (Figure~\ref{fig:seq-break}).

\begin{figure}[H]
\centering
\resizebox{\textwidth}{!}{%
\begin{sequencediagram}
  \newthread{u}{:User}
  \newinst[1.2]{home}{:Home Space}
  \newinst[1.2]{store}{:Session store}
  \newinst[1.2]{player}{:Music player}
  \newinst[1.2]{kv}{:Device storage}
  \newinst[1.2]{api}{:Backend API}

  \begin{call}{u}{tap radio}{home}{Radio screen}
    \begin{call}{home}{startOrContinue(MUSIC)}{store}{}
      \begin{call}{store}{save(start, endAt)}{kv}{}
      \end{call}
    \end{call}
  \end{call}
  \mess{u}{select station}{player}
  \begin{sdblock}{Loop}{every second}
    \begin{callself}{store}{tick(): remaining = endAt - now}{}
    \end{callself}
  \end{sdblock}
  \begin{sdblock}{Alt}{remaining = 0}
    \begin{call}{store}{stop()}{player}{}
    \end{call}
    \begin{call}{store}{enqueue(record)}{kv}{}
    \end{call}
    \begin{call}{store}{POST /api/v1/sessions}{api}{201 Created}
    \end{call}
    \mess{store}{show Session Summary}{u}
  \end{sdblock}
\end{sequencediagram}}
\caption{Sequence diagram: start a break and end it automatically}
\label{fig:seq-break}
\end{figure}

If the upload fails, the record stays in the queue and is sent again on the
next start-up or sign-in. The server ignores a second upload with the same
client ID, so a record is never stored twice.

\subsection{Sign In and Token Refresh}
\label{sec:seq-auth}

\begin{figure}[H]
\centering
\resizebox{\textwidth}{!}{%
\begin{sequencediagram}
  \newthread{u}{:User}
  \newinst[1.4]{app}{:Sign In screen}
  \newinst[1.4]{client}{:API client}
  \newinst[1.4]{api}{:AuthController}
  \newinst[1.4]{db}{:PostgreSQL}

  \begin{call}{u}{submit email, password}{app}{Home Space}
    \begin{call}{app}{login()}{client}{}
      \begin{call}{client}{POST /auth/login}{api}{access + refresh token}
        \begin{call}{api}{find user, check BCrypt hash}{db}{}
        \end{call}
        \begin{call}{api}{store SHA-256(refresh token)}{db}{}
        \end{call}
      \end{call}
    \end{call}
    \begin{callself}{app}{save refresh token (secure store)}{}
    \end{callself}
  \end{call}
  \begin{sdblock}{Opt}{a later request returns 401}
    \begin{call}{client}{POST /auth/refresh}{api}{new token pair}
      \begin{call}{api}{revoke old, store new hash}{db}{}
      \end{call}
    \end{call}
    \mess{client}{repeat original request}{api}
  \end{sdblock}
\end{sequencediagram}}
\caption{Sequence diagram: sign in and refresh the access token}
\label{fig:seq-auth}
\end{figure}

Only one refresh runs at a time: if several requests fail with 401 together,
they all wait for the same refresh. When the refresh token itself is rejected,
the app deletes the stored tokens and continues in guest mode.

\subsection{Watch Videos with a Limit}
\label{sec:seq-video}

YouTube changes pages without a full page load, so the in-app browser cannot
see video changes through normal navigation events. A small script injected
into the page reads the current address after each change and posts it to the
app; the TV screen counts each distinct video ID (Figure~\ref{fig:seq-video}).

\begin{figure}[H]
\centering
\resizebox{\textwidth}{!}{%
\begin{sequencediagram}
  \newthread{u}{:User}
  \newinst[1.3]{tv}{:TV screen}
  \newinst[1.3]{focus}{:Media focus}
  \newinst[1.3]{web}{:In-app browser}
  \newinst[1.3]{yt}{:YouTube}

  \begin{call}{u}{tap TV}{tv}{}
    \begin{call}{tv}{enterVideo()}{focus}{music paused}
    \end{call}
    \begin{call}{tv}{load m.youtube.com}{web}{}
      \begin{call}{web}{GET page}{yt}{}
      \end{call}
    \end{call}
  \end{call}
  \begin{sdblock}{Loop}{each video the user opens}
    \begin{call}{u}{open video}{web}{}
      \begin{call}{web}{postMessage(url)}{tv}{}
        \begin{callself}{tv}{count distinct video id}{}
        \end{callself}
      \end{call}
    \end{call}
  \end{sdblock}
  \begin{sdblock}{Alt}{count = 10 or break over}
    \begin{callself}{tv}{unmount browser; show end card}{}
    \end{callself}
  \end{sdblock}
  \begin{call}{u}{leave TV}{tv}{}
    \begin{call}{tv}{exitVideo()}{focus}{music resumed}
    \end{call}
  \end{call}
\end{sequencediagram}}
\caption{Sequence diagram: watch videos up to the limit}
\label{fig:seq-video}
\end{figure}

\subsection{Play a Daily Game}
\label{sec:seq-game}

\begin{figure}[H]
\centering
\resizebox{\textwidth}{!}{%
\begin{sequencediagram}
  \newthread{u}{:User}
  \newinst[1.3]{gs}{:Game screen}
  \newinst[1.3]{web}{:In-app browser}
  \newinst[1.3]{api}{:GameController}
  \newinst[1.3]{db}{:PostgreSQL}

  \begin{call}{u}{select Wordle}{gs}{}
    \begin{call}{gs}{load game URL}{web}{page loaded}
    \end{call}
    \begin{callself}{gs}{save local play, count game in break}{}
    \end{callself}
    \begin{call}{gs}{POST /games/\{id\}/plays (playDate)}{api}{201 first play / 200 repeat}
      \begin{call}{api}{INSERT ... ON CONFLICT DO NOTHING}{db}{}
      \end{call}
    \end{call}
  \end{call}
\end{sequencediagram}}
\caption{Sequence diagram: play a daily game}
\label{fig:seq-game}
\end{figure}

The play date is the user's local calendar date. The backend accepts only a
date within one day of the server's UTC date, which covers every time zone and
blocks plays for arbitrary dates. The unique key on user, game and date makes
the request safe to repeat.
